\section*{Bab \ref{ch:b3:supramolecular} --- Kimia Supramolekul}

\begin{solution}{exo:b3:supramolecular:1}
Ion--dipol; penumpukan $\pi$; kation--$\pi$; \omterm{def:b3:supramolecular:interactions}{ikatan halogen} (I ke N); tiga ikatan hidrogen.
\end{solution}

\begin{solution}{exo:b3:supramolecular:2}
$\Delta G^\circ = -RT\ln K = -8.314 \times 298 \times \ln(\num{1.0e4}) = \qty{-22.8}{kJ/mol}$.
\end{solution}

\begin{solution}{exo:b3:supramolecular:3}
$K_{\ce{K+}}/K_{\ce{Na+}} = 10^{1.7} = 50$; selisihnya $RT\ln 50 = \qty{9.7}{kJ/mol}$ yang menguntungkan kalium.
\end{solution}

\begin{solution}{exo:b3:supramolecular:4}
$n/(n + m) = 1/3$: satu \omterm{def:b3:supramolecular:supramolecular}{inang} untuk dua \omterm{def:b3:supramolecular:supramolecular}{tamu}, $\mathrm{HG_2}$.
\end{solution}

\begin{solution}{exo:b3:supramolecular:5}
$H_0 + G_0 + 1/K = \qty{4.667e-3}{mol/L}$; $[\mathrm{HG}] = (4.667 - \sqrt{4.667^2 - 4 \times 2.0 \times 2.0})/2 \times 10^{-3} = \qty{1.13e-3}{mol/L}$: fraksi
terikat 0.566, $\delta = 3.560 + 0.160 \times 0.566 = \qty{3.651}{ppm}$.
\end{solution}

\begin{solution}{exo:b3:supramolecular:6}
$\log K = 18.3 - 8.6 = 9.7$, $K = \num{5e9}$; $\Delta G^\circ = -RT\ln K = \qty{-55}{kJ/mol}$. Keenam ikatan Ni--N
sejenis di kedua sisi; jumlah molekul bebas naik dari empat menjadi
tujuh: perolehan entropi translasi itulah yang mendorong reaksi.
\end{solution}

\begin{solution}{exo:b3:supramolecular:7}
$\Delta G^\circ = -RT\ln(\num{2.0e5}) = \qty{-30.2}{kJ/mol}$; $T\Delta S^\circ = \Delta H^\circ - \Delta G^\circ = -40 + 30.2 = \qty{-9.8}{kJ/mol}$;
$\Delta S^\circ = \qty{-33}{J.K^{-1}.mol^{-1}}$: pengikatan yang didorong entalpi, yang sebagian dibayar kembali dengan entropi.
\end{solution}

\begin{solution}{exo:b3:supramolecular:8}
Tembaga(I) mengikat dua satuan fenantrolina bidentat dan, karena tetrahedral, mempertahankan
keduanya tegak lurus, yang satu terulir melalui ruang tempat cincin yang lain akan
menutup. Penutupan kedua rantai menjadi cincin lalu membuat keduanya saling mengunci. Sianida
berlebih banyak, yang mengikat tembaga(I) jauh lebih kuat, menyingkirkan logam itu dan menyisakan
\omterm{def:b3:supramolecular:interlocked}{katenana} bebas.
\end{solution}

\begin{solution}{exo:b3:supramolecular:9}
$k_c = \pi\Delta\nu/\sqrt2 = \pi \times 120/1.414 = \qty{267}{s^{-1}}$. Eyring:
$\Delta G^\ddagger = RT\ln(k_BT/hk_c) = 8.314 \times 300 \times \ln(\num{6.25e12}/267) = \qty{59.6}{kJ/mol}$.
\end{solution}

\begin{solution}{exo:b3:supramolecular:10}
Dengan obat padat berlebih, obat bebas tetap pada $S_0$. Maka $[\mathrm{DC}] = KS_0[\mathrm C]$ dan
$C_t = [\mathrm C] + [\mathrm{DC}] = [\mathrm C](1 + KS_0)$, sehingga $[\mathrm{DC}] = KS_0C_t/(1 + KS_0)$ dan kelarutan totalnya adalah $S_0 + [\mathrm{DC}]$. Di sini
$KS_0 = 0.050$: $S = 0.10 + 0.050 \times 10/1.050 = \qty{0.58}{mmol/L}$, hampir enam kali nilai intrinsiknya.
\end{solution}

\begin{solution}{exo:b3:supramolecular:11}
Jika $1/K \ll H_0, G_0$, $[\mathrm{HG}] \approx \bigl(H_0 + G_0 - |H_0 - G_0|\bigr)/2 = \min(H_0, G_0)$: fraksi terikat adalah $G_0/H_0$ sampai satu
ekuivalen, lalu 1. Kurva itu tidak lagi bergantung pada $K$: setiap nilai yang cukup besar
memberikan patahan tajam yang sama dalam batas derau, sehingga data hanya menunjukkan bahwa $K$ melampaui
suatu nilai minimum.
\end{solution}

\begin{solution}{exo:b3:supramolecular:12}
$k_{\mathrm{off}} = k_{\mathrm{on}}/K = \qty{1e9}{L.mol^{-1}.s^{-1}}/\qty{1500}{L.mol^{-1}} = \qty{6.7e5}{s^{-1}}$. Selisih geseran
$\qty{0.160}{ppm}$ adalah \qty{64}{Hz} pada \qty{400}{MHz}; koalesensi akan memerlukan $k \approx \pi \times 64/\sqrt2 = \qty{140}{s^{-1}}$. Pertukarannya
ribuan kali lebih cepat: satu sinyal rata-rata.
\end{solution}

\begin{solution}{pb:b3:supramolecular:1}
\textbf{1.} Cincin dari enam satuan $\ce{-CH2CH2O-}$; dalam kompleksnya, $\ce{K+}$ duduk di pusat, diikat oleh
keenam oksigen yang mengarah ke dalam, sedangkan gugus-gugus $\ce{CH2}$ berada di luar.

\textbf{2.} $\ce{K+}$: $-RT\ln 10 \times 6.1 = \qty{-34.8}{kJ/mol}$; $\ce{Na+}$: \qty{-25.1}{kJ/mol}.

\textbf{3.} $10^{6.1 - 4.4} = 50$.

\textbf{4.} $\ce{K+}$ yang lebih besar (\qty{138}{pm}) mencapai keenam oksigen cincin sekaligus;
$\ce{Na+}$ yang lebih kecil (\qty{102}{pm}) tidak dapat, kecuali jika cincinnya terlipat, yang memerlukan tegangan.

\textbf{5.} Bagian luarnya berupa permukaan hidrokarbon dan muatannya terbenam; $\ce{KCl}$
memerlukan ion-ionnya tersolvasi, yang tidak dapat dilakukan toluena.

\textbf{6.} Air menyolvasi $\ce{K+}$ jauh lebih baik daripada metanol, dan ion itu harus melepaskan
air itu untuk memasuki cincin.

\textbf{7.} $\ce{K+}(\text{aq}) + \ce{MnO4-}(\text{aq}) + \mathrm L(\text{org}) \rightleftharpoons \ce{[KL]+MnO4-}(\text{org})$.

\textbf{8.} $[\ce{K+}]_{\mathrm{aq}}$ tetap \qty{0.10}{mol/L}: $y = \num{1.0e5} \times 0.10 \times (\num{1.0e-3} - y)^2 = \num{1.0e4}(\num{1.0e-3} - y)^2$.

\textbf{9.} Dengan $u = \num{1.0e-3} - y$: $\num{1.0e4}u^2 + u - \num{1.0e-3} = 0$, $u = \qty{2.70e-4}{mol/L}$, $y = \qty{7.30e-4}{mol/L}$: 73\,\%
terekstraksi.

\textbf{10.} $y = \num{1.0e4}(\num{1.0e-3} - y)(\num{1.0e-2} - y)$ memberikan $y = \qty{9.89e-4}{mol/L}$: 99\,\%.

\textbf{11.} $y = \num{1.0e4}(\num{1.0e-3} - y)(\num{1.0e-4} - y)$ memberikan $y = \qty{9.0e-5}{mol/L}$: 9\,\%, dibatasi oleh \omterm{def:b3:supramolecular:hosts}{eter mahkota},
yang 90\,\% termuati.

\textbf{12.} Ion permanganat mempertahankan warna ungunya di dalam toluena. Di sana
ion itu tidak tersolvasi, sebuah anion telanjang, dan bereaksi jauh lebih cepat daripada di dalam air, tempat
selubung hidrasi melindunginya.

\textbf{13.} Ion natrium datang dan pergi jauh lebih cepat daripada selisih
kedua frekuensi resonansi: pertukaran cepat.

\textbf{14.} $\delta = \delta_{\mathrm{free}} + (\delta_{\mathrm{bound}} - \delta_{\mathrm{free}})[\mathrm{HG}]/H_0$, dengan $[\mathrm{HG}]$ dari isoterm eksak, $H_0 = \qty{2.0e-3}{mol/L}$,
$G_0 = (\text{ekuiv.})\,H_0$, $\delta_{\mathrm{free}} = \qty{3.560}{ppm}$.

\textbf{15.} $K = (1.55 \pm 0.08) \times 10^3\ \unit{L.mol^{-1}}$, $\delta_{\mathrm{bound}} = \qty{3.719 \pm 0.001}{ppm}$.

\textbf{16.} Sisaan akar rata-rata kuadratnya \qty{0.0014}{ppm}, seukuran presisi
pembacaan, tanpa kecenderungan: model 1:1 itu cocok.

\textbf{17.} Fraksi terikat $f$ pada $G_0 = fH_0 + f/(K(1 - f))$: 0.28 ekuivalen untuk $f = 0.2$ dan 2.1 untuk $f = 0.8$.

\textbf{18.} $KH_0 = 10^{6.1} \times \num{2.0e-3} \approx 2500$: kurvanya akan patah tajam pada satu ekuivalen
dan hanya memberikan batas bawah; diperlukan percobaan kompetisi atau kalorimetri pada
konsentrasi yang lebih rendah.

\textbf{19.} Di dalam toluena, yaitu larutan homogen alkena dan permanganat.

\textbf{20.} Mangan yang tereduksi meninggalkan fase organik (sebagai mangan dioksida),
dan \omterm{def:b3:supramolecular:hosts}{eter mahkota}, yang terbebas di antarmuka, mengambil $\ce{K+}$ dan $\ce{MnO4-}$ yang lain. Setiap molekul
\omterm{def:b3:supramolecular:hosts}{eter mahkota} membawa permanganat berkali-kali: sejumlah katalitik sudah cukup.

\textbf{21.} $K_{\mathrm{ex}}$ dikalikan dengan $[\ce{K+}]$: konsentrasi kalium yang tinggi mendorong ekstraksi.

\textbf{22.} Garam amonium kuaterner, seperti tetrabutilamonium klorida atau
benziltrietilamonium klorida, yang kationnya sudah lipofilik dengan sendirinya.

\textbf{23.} Benzena menyebabkan leukemia; toluena melakukan tugas yang sama dengan bahaya yang
jauh lebih kecil.

\textbf{24.} Dengan \qty{1.0}{mmol/L} \omterm{def:b3:supramolecular:hosts}{eter mahkota}, 73\,\% permanganat terekstraksi ke dalam
toluena.
\end{solution}
